\section{Прикладные применения и анализ конкретных случаев}
\subsection{Сжатие в профессиональных медиа и телекоммуникациях}
Нейросетевое сжатие находит первое практическое применение в областях, где стоимость хранения и передачи данных высока, а требования к качеству специфичны.

\begin{itemize}
\item Стриминг и видео-по-запросу (VOD): Ключевой задачей является снижение битрейта при сохранении субъективного качества. Гибридные подходы, где I-кадры (ключевые кадры) сжимаются нейросетевым кодеком (например, на базе VAE с гиперприоритетом), а P/B-кадры — традиционным видеокодеком (HEVC, VVC), показывают снижение битрейта на 20–40% для сравнимого визуального качества. Ключевым вызовом остаётся скорость декодирования для live-стримов.
\item Архивация медиафондов: Для долгосрочного хранения кино- и видеоконтента в архивах студий важна максимальная сохранность деталей. Здесь нейросетевые методы, обученные на специфичных данных (например, на оцифрованной киноплёнке), могут превзойти JPEG2000, обеспечивая лучшее визуальное качество при том же размере файла, что критично для оцифровки культурного наследия.
\end{itemize}

\subsection{Специализированные домены: медицина и дистанционное зондирование}
В этих областях изображения имеют уникальную статистику, а искажения могут привести к серьёзным ошибкам в интерпретации.

\begin{itemize}
\item Медицинская визуализация (МРТ, КТ, гистология): Традиционные кодеки, такие как JPEG, могут вносить артефакты, критичные для диагностики (размытие контуров опухолей, искажение текстур). Нейросетевые компрессоры, обученные на наборах медицинских снимков, демонстрируют лучшие результаты. Активно развивается направление \textit{diagnostically lossless compression}, где сжатие оптимизируется не для PSNR, а для минимизации ошибок последующего автоматического анализа (сегментации, классификации).
\item Спутниковые и аэрофотоснимки: Изображения имеют высокую битность (16 бит на канал) и специфичные текстуры. Нейросетевые модели могут эффективно сжимать такие данные, обеспечивая более высокое качество для последующих задач картографии, мониторинга сельского хозяйства или ЧС по сравнению с JPEG2000.
\end{itemize}

\subsection{Кейс: Сравнение в задаче сжатия портретных изображений}
Для иллюстрации проведём сравнительный анализ методов на узкой, но важной задаче — сжатии портретов, где критично сохранение деталей лица и кожи.

\begin{table}[h]
\centering
\small
\caption{Сравнение методов сжатия портретных изображений (размер 1024x1024)}
\label{tab:portrait_case}
\begin{tabular}{|l|c|c|c|}
\hline
\textbf{Метод / Параметр} & \textbf{Битрейт (bpp)} & \textbf{PSNR (дБ)} & \textbf{Субъективная оценка (1-5)} \\
\hline
JPEG (quality=90) & 0.45 & 38.2 & 3.5 (заметны артефакты) \\
\hline
WebP (quality=90) & 0.38 & 39.1 & 4.0 (лучше детали) \\
\hline
BPG (quality=medium) & 0.30 & 40.5 & 4.2 (очень хорошо) \\
\hline
\textit{Scale Hyperprior (наша, $\lambda=0.02$)} & \textit{0.28} & \textit{41.8} & \textit{4.0 (чётко, но немного "пластично")} \\
\hline
\textit{Модель с GAN (HiFiC, наша)} & \textit{0.25} & \textit{36.5} & \textit{4.7 (естественная кожа, лучшая визуально)} \\
\hline
\end{tabular}
\end{table}

\noindent \textbf{Анализ кейса:}
\begin{itemize}
\item Традиционные кодеки (JPEG, WebP) обеспечивают приемлемое качество, но на их пределе сжатия появляются артефакты (блочность, размытие ресниц).
\item Нейросетевая модель \textbf{Scale Hyperprior} показывает лучший PSNR при более низком битрейте, но может чрезмерно сглаживать текстуру кожи, делая её неестественной.
\item \textbf{GAN-модель (HiFiC)}, жертвуя объективной метрикой PSNR, достигает наивысшего субъективного качества, генерируя реалистичные текстуры, что наиболее важно для финального потребителя.
\end{itemize}

\subsection{Интеграция с существующими инфраструктурами и гибридные системы}
Полная замена существующих стандартов (JPEG, HEVC) в краткосрочной перспективе маловероятна. Более реалистичным сценарием является гибридизация.

\begin{itemize}
\item \textbf{Адаптивные гибридные кодекы:} Интеллектуальные системы, анализирующие содержимое кадра (тип сцены, наличие текста, лиц), могут динамически выбирать оптимальный метод сжатия: нейросетевой — для сложных текстур и лиц, традиционный — для простого фона или графики.
\item \textbf{Нейросетевое улучшение (upscaling) после традиционного сжатия:} Альтернативный подход — агрессивно сжимать изображение стандартным кодеком (низкий битрейт), а на стороне декодера использовать легковесную нейросеть для устранения артефактов (деблокирования, повышения резкости). Это снижает требования к скорости нейросетевой части.
\item \textbf{Краудсорсинговые и персонализированные кодеки:} Появление платформ, где пользователи могут дообучать или тонко настраивать общие модели сжатия под свой контент (например, под конкретный стиль фотографии), открывает путь к созданию персональных кодеков с максимальной эффективностью для узкого класса изображений.
\end{itemize}

Таким образом, прикладное внедрение нейросетевого сжатия идёт не по пути создания «универсального убийцы JPEG», а через развитие нишевых решений для специализированных задач (медицина, архивация) и гибридных систем, интеллектуально комбинирующих сильные стороны классических и нейросетевых подходов в зависимости от контента и условий. Успех в каждом конкретном кейсе определяется правильным выбором архитектуры, оптимизированной под целевую метрику качества — будь то объективный PSNR, диагностическая точность или субъективное визуальное восприятие.